% ============================================================
%  schematic_report.tex
%  Estructura (esqueleto) del reporte de Open RAN y RIC.
%
%  Este archivo compila POR SI SOLO en Overleaf.
%  Para incluirlo dentro de otro documento (p. ej. un main.tex
%  con % ============================================================
%  schematic_report.tex
%  Estructura (esqueleto) del reporte de Open RAN y RIC.
%
%  Este archivo compila POR SI SOLO en Overleaf.
%  Para incluirlo dentro de otro documento (p. ej. un main.tex
%  con % ============================================================
%  schematic_report.tex
%  Estructura (esqueleto) del reporte de Open RAN y RIC.
%
%  Este archivo compila POR SI SOLO en Overleaf.
%  Para incluirlo dentro de otro documento (p. ej. un main.tex
%  con % ============================================================
%  schematic_report.tex
%  Estructura (esqueleto) del reporte de Open RAN y RIC.
%
%  Este archivo compila POR SI SOLO en Overleaf.
%  Para incluirlo dentro de otro documento (p. ej. un main.tex
%  con \input{schematic_report}): borra el preambulo, es decir,
%  todo lo que va desde \documentclass hasta la linea que abre
%  el cuerpo del documento (sin incluirla), y la ultima linea.
% ============================================================

\documentclass[11pt,a4paper]{article}

\usepackage[utf8]{inputenc}   % tildes y 'n' con virgulilla (contenido en espanol)
\usepackage[T1]{fontenc}      % separacion silabica correcta en palabras acentuadas
\usepackage[spanish]{babel}   % idioma del documento

% Deja que LaTeX estire un poco mas las lineas antes de resignarse a
% desbordar el margen. Elimina los avisos "Overfull \hbox" de dos
% titulos largos. Solo actua cuando no hay otra solucion.
\setlength{\emergencystretch}{2em}

% --- Saltos de pagina entre titulos consecutivos ---
% LaTeX solo inserta la penalizacion que permite saltar de pagina ANTES
% de un titulo si \@nobreak es falso, y \@nobreak solo se restablece
% cuando arranca un parrafo. En un documento que (todavia) no tiene
% texto no arranca ningun parrafo: el flag se queda en "true" y el
% documento entero se vuelve un bloque indivisible, lo que produce una
% pagina en blanco al inicio y todo el texto desbordando la hoja.
% Restablecemos el flag antes de cada titulo. En cuanto haya texto el
% resultado es exactamente el mismo, porque \everypar ya lo hacia.
\makeatletter
\let\ORIGsection\section
\def\section{\@nobreakfalse\ORIGsection}
\let\ORIGsubsection\subsection
\def\subsection{\@nobreakfalse\ORIGsubsection}
\let\ORIGsubsubsection\subsubsection
\def\subsubsection{\@nobreakfalse\ORIGsubsubsection}
\makeatother

\begin{document}

\section{Introducción}

\subsection{Contexto y evolución de las redes de acceso radio}

\subsubsection{Evolución de D-RAN, C-RAN y NG-RAN}
\subsubsection{Limitaciones de las arquitecturas RAN tradicionales}
\subsubsection{Necesidad de desagregación, virtualización y apertura}

\subsection{Open RAN como paradigma de evolución de la RAN}

\subsubsection{Principios fundamentales de Open RAN}
\subsubsection{Desagregación e interoperabilidad multi-vendor}
\subsubsection{Softwarización, virtualización y cloudificación}

\subsection{Motivación y objetivos del reporte}

\subsubsection{Importancia de la inteligencia en Open RAN}
\subsubsection{Importancia de la seguridad en Open RAN}
\subsubsection{Alcance y organización del reporte}


\section{Arquitectura Open RAN}

\subsection{Componentes fundamentales de la arquitectura}

\subsubsection{O-RAN Radio Unit (O-RU)}
\subsubsection{O-RAN Distributed Unit (O-DU)}
\subsubsection{O-RAN Centralized Unit (O-CU)}

\subsection{Control, gestión y orquestación}

\subsubsection{Service Management and Orchestration (SMO)}
\subsubsection{O-Cloud e infraestructura virtualizada}
\subsubsection{Integración entre RAN, cloud y edge}

\subsection{Interfaces abiertas de O-RAN}

\subsubsection{Open Fronthaul}
\subsubsection{E2, A1 y O1}
\subsubsection{O2 e interacción con O-Cloud}

\subsection{Despliegue y operación de O-RAN}

\subsubsection{Distribución de componentes entre cloud, edge y sitios celulares}
\subsubsection{Arquitecturas multi-vendor}
\subsubsection{Interoperabilidad y pruebas de conformidad}


\section{RAN Intelligent Controller (RIC)}

\subsection{Concepto y función del RIC}

\subsubsection{Motivación para introducir control inteligente en la RAN}
\subsubsection{Control basado en datos y closed-loop automation}
\subsubsection{Relación entre RIC y componentes de la RAN}

\subsection{Near-Real-Time RIC}

\subsubsection{Arquitectura y componentes internos}
\subsubsection{Interfaz E2 y control de los E2 Nodes}
\subsubsection{xApps y ejecución de lógica de control}
\subsubsection{Gestión y mitigación de conflictos entre xApps}

\subsection{Non-Real-Time RIC}

\subsubsection{Arquitectura y relación con el SMO}
\subsubsection{rApps y políticas de control}
\subsubsection{Gestión de modelos y algoritmos AI/ML}

\subsection{Interacción entre Near-RT RIC y Non-RT RIC}

\subsubsection{Interfaz A1 y transferencia de políticas}
\subsubsection{Flujos de información y control}
\subsubsection{Coordinación entre diferentes escalas temporales}


\section{Aplicaciones y casos de uso del RIC}

\subsection{Gestión y optimización de recursos radio}

\subsubsection{Radio Resource Management}
\subsubsection{Load balancing}
\subsubsection{Mobility management y handover}

\subsection{Optimización de tráfico y calidad de servicio}

\subsubsection{Traffic steering}
\subsubsection{QoS y QoE management}
\subsubsection{SLA assurance}

\subsection{Optimización energética}

\subsubsection{Energy-efficient RAN}
\subsubsection{Control dinámico de recursos y celdas}
\subsubsection{Uso de AI/ML para optimización energética}

\subsection{Network slicing y servicios avanzados}

\subsubsection{Slice-aware RAN control}
\subsubsection{Orquestación y assurance de slices}
\subsubsection{Aplicaciones para V2X e Industry 5.0}

\subsection{Otros casos de uso y aplicaciones emergentes}

\subsubsection{Dynamic spectrum sharing}
\subsubsection{Detección de anomalías y optimización basada en ML}
\subsubsection{Aplicaciones para redes privadas y escenarios heterogéneos}


\section{Implementación del RIC y ecosistema de desarrollo}

\subsection{Implementaciones y plataformas de RIC}

\subsubsection{O-RAN Software Community (OSC)}
\subsubsection{ONF SD-RAN}
\subsubsection{FlexRIC y otras implementaciones}

\subsection{Desarrollo y ciclo de vida de xApps}

\subsubsection{Arquitectura de una xApp}
\subsubsection{Interfaces, APIs y modelos de servicio}
\subsubsection{Despliegue, gestión y depuración de xApps}

\subsection{AI/ML aplicado al RIC}

\subsubsection{Flujo de datos y entrenamiento de modelos}
\subsubsection{Inferencia y toma de decisiones}
\subsubsection{Closed-loop control basado en AI/ML}

\subsection{Plataformas experimentales y testbeds}

\subsubsection{OpenAirInterface y srsRAN}
\subsubsection{Colosseum y Arena}
\subsubsection{Plataformas PAWR y 5Genesis}


\section{Seguridad en Open RAN}

\subsection{Superficie de ataque de Open RAN}

\subsubsection{Impacto de la desagregación y multi-vendorismo}
\subsubsection{Virtualización, cloudificación y software de terceros}
\subsubsection{Nuevos vectores de ataque introducidos por las interfaces abiertas}

\subsection{Amenazas sobre los componentes e interfaces}

\subsubsection{Amenazas sobre O-RU, O-DU y O-CU}
\subsubsection{Amenazas sobre las interfaces E2, A1, O1 y O2}
\subsubsection{Amenazas sobre el Open Fronthaul}

\subsection{Amenazas asociadas al RIC y sus aplicaciones}

\subsubsection{xApps y rApps maliciosas o comprometidas}
\subsubsection{Ataques contra E2 Nodes y señalización}
\subsubsection{Conflictos y manipulación de decisiones de control}

\subsection{Ataques y riesgos de seguridad}

\subsubsection{Denial of Service y Distributed Denial of Service}
\subsubsection{Man-in-the-Middle y ataques sobre las comunicaciones}
\subsubsection{Spoofing, inyección de datos y manipulación de tráfico}
\subsubsection{Ataques relacionados con AI/ML}

\subsection{Privacidad y protección de información}

\subsubsection{Datos de usuarios y telemetría de la RAN}
\subsubsection{Riesgos derivados de la recopilación y exposición de datos}
\subsubsection{Privacidad en escenarios de inteligencia distribuida}


\section{Mecanismos y estrategias de seguridad para Open RAN}

\subsection{Seguridad de interfaces y comunicaciones}

\subsubsection{Autenticación y autorización}
\subsubsection{Protección de las comunicaciones mediante mecanismos criptográficos}
\subsubsection{Integridad y confidencialidad de la información}

\subsection{Control de acceso y confianza en el ecosistema O-RAN}

\subsubsection{Gestión de identidades y privilegios}
\subsubsection{Zero Trust aplicado a Open RAN}
\subsubsection{Verificación y aislamiento de aplicaciones}

\subsection{Seguridad de xApps y rApps}

\subsubsection{Validación y autenticación de aplicaciones}
\subsubsection{Gestión de permisos y privilegios}
\subsubsection{Detección de comportamiento malicioso}

\subsection{Detección y mitigación de ataques}

\subsubsection{Monitoreo de interfaces y tráfico}
\subsubsection{Detección basada en anomalías}
\subsubsection{Uso de AI/ML para seguridad}

\subsection{Estandarización y buenas prácticas}

\subsubsection{Esfuerzos de estandarización de seguridad}
\subsubsection{Interoperabilidad y pruebas de seguridad}
\subsubsection{Gestión del ciclo de vida y actualización de componentes}


\section{Relación entre RIC, inteligencia y seguridad}

\subsection{El RIC como habilitador de mecanismos de seguridad}

\subsubsection{Monitoreo inteligente de la RAN}
\subsubsection{Detección de anomalías en tiempo casi real}
\subsubsection{Respuesta y mitigación mediante closed-loop control}

\subsection{Seguridad de los mecanismos de inteligencia}

\subsubsection{Ataques contra modelos y datos AI/ML}
\subsubsection{Manipulación de modelos y decisiones}
\subsubsection{Seguridad del ciclo de vida AI/ML}

\subsection{Trade-offs entre inteligencia, apertura y seguridad}

\subsubsection{Beneficios de la apertura frente a la ampliación de la superficie de ataque}
\subsubsection{Flexibilidad multi-vendor frente a complejidad de integración}
\subsubsection{Automatización frente a riesgos asociados al control autónomo}


\section{Desafíos y oportunidades de investigación}

\subsection{Desafíos arquitectónicos y de interoperabilidad}

\subsubsection{Interoperabilidad multi-vendor}
\subsubsection{Evolución de interfaces y E2 Service Models}
\subsubsection{Escalabilidad y distribución cloud-edge}

\subsection{Desafíos de inteligencia y control}

\subsubsection{Latencia y restricciones temporales}
\subsubsection{Conflictos entre aplicaciones del RIC}
\subsubsection{Generalización y confiabilidad de modelos AI/ML}

\subsection{Desafíos de seguridad}

\subsubsection{Protección de componentes y aplicaciones de terceros}
\subsubsection{Seguridad de interfaces abiertas}
\subsubsection{Ataques contra AI/ML y control inteligente}

\subsection{Direcciones futuras}

\subsubsection{Evolución hacia redes 6G}
\subsubsection{Nuevos casos de uso y arquitecturas}
\subsubsection{Integración de inteligencia y seguridad por diseño}


\section{Análisis y discusión}

\subsection{Síntesis de la evolución hacia O-RAN inteligente y seguro}

\subsubsection{Relación entre apertura, desagregación e inteligencia}
\subsubsection{Relación entre RIC, automatización y optimización}
\subsubsection{Relación entre apertura, inteligencia y seguridad}

\subsection{Comparación de enfoques de implementación}

\subsubsection{Características de las principales implementaciones}
\subsubsection{Capacidades para desarrollo y experimentación}
\subsubsection{Limitaciones identificadas en la literatura}

\subsection{Análisis de los principales desafíos}

\subsubsection{Desafíos técnicos}
\subsubsection{Desafíos de seguridad}
\subsubsection{Desafíos de estandarización e interoperabilidad}


\section{Conclusiones}

\subsection{Principales hallazgos}

\subsubsection{Características fundamentales de Open RAN}
\subsubsection{Rol del RIC en la evolución de la RAN}
\subsubsection{Principales implicaciones de seguridad}

\subsection{Perspectivas futuras}

\subsubsection{Evolución del RIC y las arquitecturas O-RAN}
\subsubsection{Evolución de las estrategias de seguridad}
\subsubsection{Oportunidades de investigación en 5G-Advanced y 6G}

\end{document}
): borra el preambulo, es decir,
%  todo lo que va desde \documentclass hasta la linea que abre
%  el cuerpo del documento (sin incluirla), y la ultima linea.
% ============================================================

\documentclass[11pt,a4paper]{article}

\usepackage[utf8]{inputenc}   % tildes y 'n' con virgulilla (contenido en espanol)
\usepackage[T1]{fontenc}      % separacion silabica correcta en palabras acentuadas
\usepackage[spanish]{babel}   % idioma del documento

% Deja que LaTeX estire un poco mas las lineas antes de resignarse a
% desbordar el margen. Elimina los avisos "Overfull \hbox" de dos
% titulos largos. Solo actua cuando no hay otra solucion.
\setlength{\emergencystretch}{2em}

% --- Saltos de pagina entre titulos consecutivos ---
% LaTeX solo inserta la penalizacion que permite saltar de pagina ANTES
% de un titulo si \@nobreak es falso, y \@nobreak solo se restablece
% cuando arranca un parrafo. En un documento que (todavia) no tiene
% texto no arranca ningun parrafo: el flag se queda en "true" y el
% documento entero se vuelve un bloque indivisible, lo que produce una
% pagina en blanco al inicio y todo el texto desbordando la hoja.
% Restablecemos el flag antes de cada titulo. En cuanto haya texto el
% resultado es exactamente el mismo, porque \everypar ya lo hacia.
\makeatletter
\let\ORIGsection\section
\def\section{\@nobreakfalse\ORIGsection}
\let\ORIGsubsection\subsection
\def\subsection{\@nobreakfalse\ORIGsubsection}
\let\ORIGsubsubsection\subsubsection
\def\subsubsection{\@nobreakfalse\ORIGsubsubsection}
\makeatother

\begin{document}

\section{Introducción}

\subsection{Contexto y evolución de las redes de acceso radio}

\subsubsection{Evolución de D-RAN, C-RAN y NG-RAN}
\subsubsection{Limitaciones de las arquitecturas RAN tradicionales}
\subsubsection{Necesidad de desagregación, virtualización y apertura}

\subsection{Open RAN como paradigma de evolución de la RAN}

\subsubsection{Principios fundamentales de Open RAN}
\subsubsection{Desagregación e interoperabilidad multi-vendor}
\subsubsection{Softwarización, virtualización y cloudificación}

\subsection{Motivación y objetivos del reporte}

\subsubsection{Importancia de la inteligencia en Open RAN}
\subsubsection{Importancia de la seguridad en Open RAN}
\subsubsection{Alcance y organización del reporte}


\section{Arquitectura Open RAN}

\subsection{Componentes fundamentales de la arquitectura}

\subsubsection{O-RAN Radio Unit (O-RU)}
\subsubsection{O-RAN Distributed Unit (O-DU)}
\subsubsection{O-RAN Centralized Unit (O-CU)}

\subsection{Control, gestión y orquestación}

\subsubsection{Service Management and Orchestration (SMO)}
\subsubsection{O-Cloud e infraestructura virtualizada}
\subsubsection{Integración entre RAN, cloud y edge}

\subsection{Interfaces abiertas de O-RAN}

\subsubsection{Open Fronthaul}
\subsubsection{E2, A1 y O1}
\subsubsection{O2 e interacción con O-Cloud}

\subsection{Despliegue y operación de O-RAN}

\subsubsection{Distribución de componentes entre cloud, edge y sitios celulares}
\subsubsection{Arquitecturas multi-vendor}
\subsubsection{Interoperabilidad y pruebas de conformidad}


\section{RAN Intelligent Controller (RIC)}

\subsection{Concepto y función del RIC}

\subsubsection{Motivación para introducir control inteligente en la RAN}
\subsubsection{Control basado en datos y closed-loop automation}
\subsubsection{Relación entre RIC y componentes de la RAN}

\subsection{Near-Real-Time RIC}

\subsubsection{Arquitectura y componentes internos}
\subsubsection{Interfaz E2 y control de los E2 Nodes}
\subsubsection{xApps y ejecución de lógica de control}
\subsubsection{Gestión y mitigación de conflictos entre xApps}

\subsection{Non-Real-Time RIC}

\subsubsection{Arquitectura y relación con el SMO}
\subsubsection{rApps y políticas de control}
\subsubsection{Gestión de modelos y algoritmos AI/ML}

\subsection{Interacción entre Near-RT RIC y Non-RT RIC}

\subsubsection{Interfaz A1 y transferencia de políticas}
\subsubsection{Flujos de información y control}
\subsubsection{Coordinación entre diferentes escalas temporales}


\section{Aplicaciones y casos de uso del RIC}

\subsection{Gestión y optimización de recursos radio}

\subsubsection{Radio Resource Management}
\subsubsection{Load balancing}
\subsubsection{Mobility management y handover}

\subsection{Optimización de tráfico y calidad de servicio}

\subsubsection{Traffic steering}
\subsubsection{QoS y QoE management}
\subsubsection{SLA assurance}

\subsection{Optimización energética}

\subsubsection{Energy-efficient RAN}
\subsubsection{Control dinámico de recursos y celdas}
\subsubsection{Uso de AI/ML para optimización energética}

\subsection{Network slicing y servicios avanzados}

\subsubsection{Slice-aware RAN control}
\subsubsection{Orquestación y assurance de slices}
\subsubsection{Aplicaciones para V2X e Industry 5.0}

\subsection{Otros casos de uso y aplicaciones emergentes}

\subsubsection{Dynamic spectrum sharing}
\subsubsection{Detección de anomalías y optimización basada en ML}
\subsubsection{Aplicaciones para redes privadas y escenarios heterogéneos}


\section{Implementación del RIC y ecosistema de desarrollo}

\subsection{Implementaciones y plataformas de RIC}

\subsubsection{O-RAN Software Community (OSC)}
\subsubsection{ONF SD-RAN}
\subsubsection{FlexRIC y otras implementaciones}

\subsection{Desarrollo y ciclo de vida de xApps}

\subsubsection{Arquitectura de una xApp}
\subsubsection{Interfaces, APIs y modelos de servicio}
\subsubsection{Despliegue, gestión y depuración de xApps}

\subsection{AI/ML aplicado al RIC}

\subsubsection{Flujo de datos y entrenamiento de modelos}
\subsubsection{Inferencia y toma de decisiones}
\subsubsection{Closed-loop control basado en AI/ML}

\subsection{Plataformas experimentales y testbeds}

\subsubsection{OpenAirInterface y srsRAN}
\subsubsection{Colosseum y Arena}
\subsubsection{Plataformas PAWR y 5Genesis}


\section{Seguridad en Open RAN}

\subsection{Superficie de ataque de Open RAN}

\subsubsection{Impacto de la desagregación y multi-vendorismo}
\subsubsection{Virtualización, cloudificación y software de terceros}
\subsubsection{Nuevos vectores de ataque introducidos por las interfaces abiertas}

\subsection{Amenazas sobre los componentes e interfaces}

\subsubsection{Amenazas sobre O-RU, O-DU y O-CU}
\subsubsection{Amenazas sobre las interfaces E2, A1, O1 y O2}
\subsubsection{Amenazas sobre el Open Fronthaul}

\subsection{Amenazas asociadas al RIC y sus aplicaciones}

\subsubsection{xApps y rApps maliciosas o comprometidas}
\subsubsection{Ataques contra E2 Nodes y señalización}
\subsubsection{Conflictos y manipulación de decisiones de control}

\subsection{Ataques y riesgos de seguridad}

\subsubsection{Denial of Service y Distributed Denial of Service}
\subsubsection{Man-in-the-Middle y ataques sobre las comunicaciones}
\subsubsection{Spoofing, inyección de datos y manipulación de tráfico}
\subsubsection{Ataques relacionados con AI/ML}

\subsection{Privacidad y protección de información}

\subsubsection{Datos de usuarios y telemetría de la RAN}
\subsubsection{Riesgos derivados de la recopilación y exposición de datos}
\subsubsection{Privacidad en escenarios de inteligencia distribuida}


\section{Mecanismos y estrategias de seguridad para Open RAN}

\subsection{Seguridad de interfaces y comunicaciones}

\subsubsection{Autenticación y autorización}
\subsubsection{Protección de las comunicaciones mediante mecanismos criptográficos}
\subsubsection{Integridad y confidencialidad de la información}

\subsection{Control de acceso y confianza en el ecosistema O-RAN}

\subsubsection{Gestión de identidades y privilegios}
\subsubsection{Zero Trust aplicado a Open RAN}
\subsubsection{Verificación y aislamiento de aplicaciones}

\subsection{Seguridad de xApps y rApps}

\subsubsection{Validación y autenticación de aplicaciones}
\subsubsection{Gestión de permisos y privilegios}
\subsubsection{Detección de comportamiento malicioso}

\subsection{Detección y mitigación de ataques}

\subsubsection{Monitoreo de interfaces y tráfico}
\subsubsection{Detección basada en anomalías}
\subsubsection{Uso de AI/ML para seguridad}

\subsection{Estandarización y buenas prácticas}

\subsubsection{Esfuerzos de estandarización de seguridad}
\subsubsection{Interoperabilidad y pruebas de seguridad}
\subsubsection{Gestión del ciclo de vida y actualización de componentes}


\section{Relación entre RIC, inteligencia y seguridad}

\subsection{El RIC como habilitador de mecanismos de seguridad}

\subsubsection{Monitoreo inteligente de la RAN}
\subsubsection{Detección de anomalías en tiempo casi real}
\subsubsection{Respuesta y mitigación mediante closed-loop control}

\subsection{Seguridad de los mecanismos de inteligencia}

\subsubsection{Ataques contra modelos y datos AI/ML}
\subsubsection{Manipulación de modelos y decisiones}
\subsubsection{Seguridad del ciclo de vida AI/ML}

\subsection{Trade-offs entre inteligencia, apertura y seguridad}

\subsubsection{Beneficios de la apertura frente a la ampliación de la superficie de ataque}
\subsubsection{Flexibilidad multi-vendor frente a complejidad de integración}
\subsubsection{Automatización frente a riesgos asociados al control autónomo}


\section{Desafíos y oportunidades de investigación}

\subsection{Desafíos arquitectónicos y de interoperabilidad}

\subsubsection{Interoperabilidad multi-vendor}
\subsubsection{Evolución de interfaces y E2 Service Models}
\subsubsection{Escalabilidad y distribución cloud-edge}

\subsection{Desafíos de inteligencia y control}

\subsubsection{Latencia y restricciones temporales}
\subsubsection{Conflictos entre aplicaciones del RIC}
\subsubsection{Generalización y confiabilidad de modelos AI/ML}

\subsection{Desafíos de seguridad}

\subsubsection{Protección de componentes y aplicaciones de terceros}
\subsubsection{Seguridad de interfaces abiertas}
\subsubsection{Ataques contra AI/ML y control inteligente}

\subsection{Direcciones futuras}

\subsubsection{Evolución hacia redes 6G}
\subsubsection{Nuevos casos de uso y arquitecturas}
\subsubsection{Integración de inteligencia y seguridad por diseño}


\section{Análisis y discusión}

\subsection{Síntesis de la evolución hacia O-RAN inteligente y seguro}

\subsubsection{Relación entre apertura, desagregación e inteligencia}
\subsubsection{Relación entre RIC, automatización y optimización}
\subsubsection{Relación entre apertura, inteligencia y seguridad}

\subsection{Comparación de enfoques de implementación}

\subsubsection{Características de las principales implementaciones}
\subsubsection{Capacidades para desarrollo y experimentación}
\subsubsection{Limitaciones identificadas en la literatura}

\subsection{Análisis de los principales desafíos}

\subsubsection{Desafíos técnicos}
\subsubsection{Desafíos de seguridad}
\subsubsection{Desafíos de estandarización e interoperabilidad}


\section{Conclusiones}

\subsection{Principales hallazgos}

\subsubsection{Características fundamentales de Open RAN}
\subsubsection{Rol del RIC en la evolución de la RAN}
\subsubsection{Principales implicaciones de seguridad}

\subsection{Perspectivas futuras}

\subsubsection{Evolución del RIC y las arquitecturas O-RAN}
\subsubsection{Evolución de las estrategias de seguridad}
\subsubsection{Oportunidades de investigación en 5G-Advanced y 6G}

\end{document}
): borra el preambulo, es decir,
%  todo lo que va desde \documentclass hasta la linea que abre
%  el cuerpo del documento (sin incluirla), y la ultima linea.
% ============================================================

\documentclass[11pt,a4paper]{article}

\usepackage[utf8]{inputenc}   % tildes y 'n' con virgulilla (contenido en espanol)
\usepackage[T1]{fontenc}      % separacion silabica correcta en palabras acentuadas
\usepackage[spanish]{babel}   % idioma del documento

% Deja que LaTeX estire un poco mas las lineas antes de resignarse a
% desbordar el margen. Elimina los avisos "Overfull \hbox" de dos
% titulos largos. Solo actua cuando no hay otra solucion.
\setlength{\emergencystretch}{2em}

% --- Saltos de pagina entre titulos consecutivos ---
% LaTeX solo inserta la penalizacion que permite saltar de pagina ANTES
% de un titulo si \@nobreak es falso, y \@nobreak solo se restablece
% cuando arranca un parrafo. En un documento que (todavia) no tiene
% texto no arranca ningun parrafo: el flag se queda en "true" y el
% documento entero se vuelve un bloque indivisible, lo que produce una
% pagina en blanco al inicio y todo el texto desbordando la hoja.
% Restablecemos el flag antes de cada titulo. En cuanto haya texto el
% resultado es exactamente el mismo, porque \everypar ya lo hacia.
\makeatletter
\let\ORIGsection\section
\def\section{\@nobreakfalse\ORIGsection}
\let\ORIGsubsection\subsection
\def\subsection{\@nobreakfalse\ORIGsubsection}
\let\ORIGsubsubsection\subsubsection
\def\subsubsection{\@nobreakfalse\ORIGsubsubsection}
\makeatother

\begin{document}

\section{Introducción}

\subsection{Contexto y evolución de las redes de acceso radio}

\subsubsection{Evolución de D-RAN, C-RAN y NG-RAN}
\subsubsection{Limitaciones de las arquitecturas RAN tradicionales}
\subsubsection{Necesidad de desagregación, virtualización y apertura}

\subsection{Open RAN como paradigma de evolución de la RAN}

\subsubsection{Principios fundamentales de Open RAN}
\subsubsection{Desagregación e interoperabilidad multi-vendor}
\subsubsection{Softwarización, virtualización y cloudificación}

\subsection{Motivación y objetivos del reporte}

\subsubsection{Importancia de la inteligencia en Open RAN}
\subsubsection{Importancia de la seguridad en Open RAN}
\subsubsection{Alcance y organización del reporte}


\section{Arquitectura Open RAN}

\subsection{Componentes fundamentales de la arquitectura}

\subsubsection{O-RAN Radio Unit (O-RU)}
\subsubsection{O-RAN Distributed Unit (O-DU)}
\subsubsection{O-RAN Centralized Unit (O-CU)}

\subsection{Control, gestión y orquestación}

\subsubsection{Service Management and Orchestration (SMO)}
\subsubsection{O-Cloud e infraestructura virtualizada}
\subsubsection{Integración entre RAN, cloud y edge}

\subsection{Interfaces abiertas de O-RAN}

\subsubsection{Open Fronthaul}
\subsubsection{E2, A1 y O1}
\subsubsection{O2 e interacción con O-Cloud}

\subsection{Despliegue y operación de O-RAN}

\subsubsection{Distribución de componentes entre cloud, edge y sitios celulares}
\subsubsection{Arquitecturas multi-vendor}
\subsubsection{Interoperabilidad y pruebas de conformidad}


\section{RAN Intelligent Controller (RIC)}

\subsection{Concepto y función del RIC}

\subsubsection{Motivación para introducir control inteligente en la RAN}
\subsubsection{Control basado en datos y closed-loop automation}
\subsubsection{Relación entre RIC y componentes de la RAN}

\subsection{Near-Real-Time RIC}

\subsubsection{Arquitectura y componentes internos}
\subsubsection{Interfaz E2 y control de los E2 Nodes}
\subsubsection{xApps y ejecución de lógica de control}
\subsubsection{Gestión y mitigación de conflictos entre xApps}

\subsection{Non-Real-Time RIC}

\subsubsection{Arquitectura y relación con el SMO}
\subsubsection{rApps y políticas de control}
\subsubsection{Gestión de modelos y algoritmos AI/ML}

\subsection{Interacción entre Near-RT RIC y Non-RT RIC}

\subsubsection{Interfaz A1 y transferencia de políticas}
\subsubsection{Flujos de información y control}
\subsubsection{Coordinación entre diferentes escalas temporales}


\section{Aplicaciones y casos de uso del RIC}

\subsection{Gestión y optimización de recursos radio}

\subsubsection{Radio Resource Management}
\subsubsection{Load balancing}
\subsubsection{Mobility management y handover}

\subsection{Optimización de tráfico y calidad de servicio}

\subsubsection{Traffic steering}
\subsubsection{QoS y QoE management}
\subsubsection{SLA assurance}

\subsection{Optimización energética}

\subsubsection{Energy-efficient RAN}
\subsubsection{Control dinámico de recursos y celdas}
\subsubsection{Uso de AI/ML para optimización energética}

\subsection{Network slicing y servicios avanzados}

\subsubsection{Slice-aware RAN control}
\subsubsection{Orquestación y assurance de slices}
\subsubsection{Aplicaciones para V2X e Industry 5.0}

\subsection{Otros casos de uso y aplicaciones emergentes}

\subsubsection{Dynamic spectrum sharing}
\subsubsection{Detección de anomalías y optimización basada en ML}
\subsubsection{Aplicaciones para redes privadas y escenarios heterogéneos}


\section{Implementación del RIC y ecosistema de desarrollo}

\subsection{Implementaciones y plataformas de RIC}

\subsubsection{O-RAN Software Community (OSC)}
\subsubsection{ONF SD-RAN}
\subsubsection{FlexRIC y otras implementaciones}

\subsection{Desarrollo y ciclo de vida de xApps}

\subsubsection{Arquitectura de una xApp}
\subsubsection{Interfaces, APIs y modelos de servicio}
\subsubsection{Despliegue, gestión y depuración de xApps}

\subsection{AI/ML aplicado al RIC}

\subsubsection{Flujo de datos y entrenamiento de modelos}
\subsubsection{Inferencia y toma de decisiones}
\subsubsection{Closed-loop control basado en AI/ML}

\subsection{Plataformas experimentales y testbeds}

\subsubsection{OpenAirInterface y srsRAN}
\subsubsection{Colosseum y Arena}
\subsubsection{Plataformas PAWR y 5Genesis}


\section{Seguridad en Open RAN}

\subsection{Superficie de ataque de Open RAN}

\subsubsection{Impacto de la desagregación y multi-vendorismo}
\subsubsection{Virtualización, cloudificación y software de terceros}
\subsubsection{Nuevos vectores de ataque introducidos por las interfaces abiertas}

\subsection{Amenazas sobre los componentes e interfaces}

\subsubsection{Amenazas sobre O-RU, O-DU y O-CU}
\subsubsection{Amenazas sobre las interfaces E2, A1, O1 y O2}
\subsubsection{Amenazas sobre el Open Fronthaul}

\subsection{Amenazas asociadas al RIC y sus aplicaciones}

\subsubsection{xApps y rApps maliciosas o comprometidas}
\subsubsection{Ataques contra E2 Nodes y señalización}
\subsubsection{Conflictos y manipulación de decisiones de control}

\subsection{Ataques y riesgos de seguridad}

\subsubsection{Denial of Service y Distributed Denial of Service}
\subsubsection{Man-in-the-Middle y ataques sobre las comunicaciones}
\subsubsection{Spoofing, inyección de datos y manipulación de tráfico}
\subsubsection{Ataques relacionados con AI/ML}

\subsection{Privacidad y protección de información}

\subsubsection{Datos de usuarios y telemetría de la RAN}
\subsubsection{Riesgos derivados de la recopilación y exposición de datos}
\subsubsection{Privacidad en escenarios de inteligencia distribuida}


\section{Mecanismos y estrategias de seguridad para Open RAN}

\subsection{Seguridad de interfaces y comunicaciones}

\subsubsection{Autenticación y autorización}
\subsubsection{Protección de las comunicaciones mediante mecanismos criptográficos}
\subsubsection{Integridad y confidencialidad de la información}

\subsection{Control de acceso y confianza en el ecosistema O-RAN}

\subsubsection{Gestión de identidades y privilegios}
\subsubsection{Zero Trust aplicado a Open RAN}
\subsubsection{Verificación y aislamiento de aplicaciones}

\subsection{Seguridad de xApps y rApps}

\subsubsection{Validación y autenticación de aplicaciones}
\subsubsection{Gestión de permisos y privilegios}
\subsubsection{Detección de comportamiento malicioso}

\subsection{Detección y mitigación de ataques}

\subsubsection{Monitoreo de interfaces y tráfico}
\subsubsection{Detección basada en anomalías}
\subsubsection{Uso de AI/ML para seguridad}

\subsection{Estandarización y buenas prácticas}

\subsubsection{Esfuerzos de estandarización de seguridad}
\subsubsection{Interoperabilidad y pruebas de seguridad}
\subsubsection{Gestión del ciclo de vida y actualización de componentes}


\section{Relación entre RIC, inteligencia y seguridad}

\subsection{El RIC como habilitador de mecanismos de seguridad}

\subsubsection{Monitoreo inteligente de la RAN}
\subsubsection{Detección de anomalías en tiempo casi real}
\subsubsection{Respuesta y mitigación mediante closed-loop control}

\subsection{Seguridad de los mecanismos de inteligencia}

\subsubsection{Ataques contra modelos y datos AI/ML}
\subsubsection{Manipulación de modelos y decisiones}
\subsubsection{Seguridad del ciclo de vida AI/ML}

\subsection{Trade-offs entre inteligencia, apertura y seguridad}

\subsubsection{Beneficios de la apertura frente a la ampliación de la superficie de ataque}
\subsubsection{Flexibilidad multi-vendor frente a complejidad de integración}
\subsubsection{Automatización frente a riesgos asociados al control autónomo}


\section{Desafíos y oportunidades de investigación}

\subsection{Desafíos arquitectónicos y de interoperabilidad}

\subsubsection{Interoperabilidad multi-vendor}
\subsubsection{Evolución de interfaces y E2 Service Models}
\subsubsection{Escalabilidad y distribución cloud-edge}

\subsection{Desafíos de inteligencia y control}

\subsubsection{Latencia y restricciones temporales}
\subsubsection{Conflictos entre aplicaciones del RIC}
\subsubsection{Generalización y confiabilidad de modelos AI/ML}

\subsection{Desafíos de seguridad}

\subsubsection{Protección de componentes y aplicaciones de terceros}
\subsubsection{Seguridad de interfaces abiertas}
\subsubsection{Ataques contra AI/ML y control inteligente}

\subsection{Direcciones futuras}

\subsubsection{Evolución hacia redes 6G}
\subsubsection{Nuevos casos de uso y arquitecturas}
\subsubsection{Integración de inteligencia y seguridad por diseño}


\section{Análisis y discusión}

\subsection{Síntesis de la evolución hacia O-RAN inteligente y seguro}

\subsubsection{Relación entre apertura, desagregación e inteligencia}
\subsubsection{Relación entre RIC, automatización y optimización}
\subsubsection{Relación entre apertura, inteligencia y seguridad}

\subsection{Comparación de enfoques de implementación}

\subsubsection{Características de las principales implementaciones}
\subsubsection{Capacidades para desarrollo y experimentación}
\subsubsection{Limitaciones identificadas en la literatura}

\subsection{Análisis de los principales desafíos}

\subsubsection{Desafíos técnicos}
\subsubsection{Desafíos de seguridad}
\subsubsection{Desafíos de estandarización e interoperabilidad}


\section{Conclusiones}

\subsection{Principales hallazgos}

\subsubsection{Características fundamentales de Open RAN}
\subsubsection{Rol del RIC en la evolución de la RAN}
\subsubsection{Principales implicaciones de seguridad}

\subsection{Perspectivas futuras}

\subsubsection{Evolución del RIC y las arquitecturas O-RAN}
\subsubsection{Evolución de las estrategias de seguridad}
\subsubsection{Oportunidades de investigación en 5G-Advanced y 6G}

\end{document}
): borra el preambulo, es decir,
%  todo lo que va desde \documentclass hasta la linea que abre
%  el cuerpo del documento (sin incluirla), y la ultima linea.
% ============================================================

\documentclass[11pt,a4paper]{article}

\usepackage[utf8]{inputenc}   % tildes y 'n' con virgulilla (contenido en espanol)
\usepackage[T1]{fontenc}      % separacion silabica correcta en palabras acentuadas
\usepackage[spanish]{babel}   % idioma del documento

% Deja que LaTeX estire un poco mas las lineas antes de resignarse a
% desbordar el margen. Elimina los avisos "Overfull \hbox" de dos
% titulos largos. Solo actua cuando no hay otra solucion.
\setlength{\emergencystretch}{2em}

% --- Saltos de pagina entre titulos consecutivos ---
% LaTeX solo inserta la penalizacion que permite saltar de pagina ANTES
% de un titulo si \@nobreak es falso, y \@nobreak solo se restablece
% cuando arranca un parrafo. En un documento que (todavia) no tiene
% texto no arranca ningun parrafo: el flag se queda en "true" y el
% documento entero se vuelve un bloque indivisible, lo que produce una
% pagina en blanco al inicio y todo el texto desbordando la hoja.
% Restablecemos el flag antes de cada titulo. En cuanto haya texto el
% resultado es exactamente el mismo, porque \everypar ya lo hacia.
\makeatletter
\let\ORIGsection\section
\def\section{\@nobreakfalse\ORIGsection}
\let\ORIGsubsection\subsection
\def\subsection{\@nobreakfalse\ORIGsubsection}
\let\ORIGsubsubsection\subsubsection
\def\subsubsection{\@nobreakfalse\ORIGsubsubsection}
\makeatother

\begin{document}

\section{Introducción}

\subsection{Contexto y evolución de las redes de acceso radio}

\subsubsection{Evolución de D-RAN, C-RAN y NG-RAN}
\subsubsection{Limitaciones de las arquitecturas RAN tradicionales}
\subsubsection{Necesidad de desagregación, virtualización y apertura}

\subsection{Open RAN como paradigma de evolución de la RAN}

\subsubsection{Principios fundamentales de Open RAN}
\subsubsection{Desagregación e interoperabilidad multi-vendor}
\subsubsection{Softwarización, virtualización y cloudificación}

\subsection{Motivación y objetivos del reporte}

\subsubsection{Importancia de la inteligencia en Open RAN}
\subsubsection{Importancia de la seguridad en Open RAN}
\subsubsection{Alcance y organización del reporte}


\section{Arquitectura Open RAN}

\subsection{Componentes fundamentales de la arquitectura}

\subsubsection{O-RAN Radio Unit (O-RU)}
\subsubsection{O-RAN Distributed Unit (O-DU)}
\subsubsection{O-RAN Centralized Unit (O-CU)}

\subsection{Control, gestión y orquestación}

\subsubsection{Service Management and Orchestration (SMO)}
\subsubsection{O-Cloud e infraestructura virtualizada}
\subsubsection{Integración entre RAN, cloud y edge}

\subsection{Interfaces abiertas de O-RAN}

\subsubsection{Open Fronthaul}
\subsubsection{E2, A1 y O1}
\subsubsection{O2 e interacción con O-Cloud}

\subsection{Despliegue y operación de O-RAN}

\subsubsection{Distribución de componentes entre cloud, edge y sitios celulares}
\subsubsection{Arquitecturas multi-vendor}
\subsubsection{Interoperabilidad y pruebas de conformidad}


\section{RAN Intelligent Controller (RIC)}

\subsection{Concepto y función del RIC}

\subsubsection{Motivación para introducir control inteligente en la RAN}
\subsubsection{Control basado en datos y closed-loop automation}
\subsubsection{Relación entre RIC y componentes de la RAN}

\subsection{Near-Real-Time RIC}

\subsubsection{Arquitectura y componentes internos}
\subsubsection{Interfaz E2 y control de los E2 Nodes}
\subsubsection{xApps y ejecución de lógica de control}
\subsubsection{Gestión y mitigación de conflictos entre xApps}

\subsection{Non-Real-Time RIC}

\subsubsection{Arquitectura y relación con el SMO}
\subsubsection{rApps y políticas de control}
\subsubsection{Gestión de modelos y algoritmos AI/ML}

\subsection{Interacción entre Near-RT RIC y Non-RT RIC}

\subsubsection{Interfaz A1 y transferencia de políticas}
\subsubsection{Flujos de información y control}
\subsubsection{Coordinación entre diferentes escalas temporales}


\section{Aplicaciones y casos de uso del RIC}

\subsection{Gestión y optimización de recursos radio}

\subsubsection{Radio Resource Management}
\subsubsection{Load balancing}
\subsubsection{Mobility management y handover}

\subsection{Optimización de tráfico y calidad de servicio}

\subsubsection{Traffic steering}
\subsubsection{QoS y QoE management}
\subsubsection{SLA assurance}

\subsection{Optimización energética}

\subsubsection{Energy-efficient RAN}
\subsubsection{Control dinámico de recursos y celdas}
\subsubsection{Uso de AI/ML para optimización energética}

\subsection{Network slicing y servicios avanzados}

\subsubsection{Slice-aware RAN control}
\subsubsection{Orquestación y assurance de slices}
\subsubsection{Aplicaciones para V2X e Industry 5.0}

\subsection{Otros casos de uso y aplicaciones emergentes}

\subsubsection{Dynamic spectrum sharing}
\subsubsection{Detección de anomalías y optimización basada en ML}
\subsubsection{Aplicaciones para redes privadas y escenarios heterogéneos}


\section{Implementación del RIC y ecosistema de desarrollo}

\subsection{Implementaciones y plataformas de RIC}

\subsubsection{O-RAN Software Community (OSC)}
\subsubsection{ONF SD-RAN}
\subsubsection{FlexRIC y otras implementaciones}

\subsection{Desarrollo y ciclo de vida de xApps}

\subsubsection{Arquitectura de una xApp}
\subsubsection{Interfaces, APIs y modelos de servicio}
\subsubsection{Despliegue, gestión y depuración de xApps}

\subsection{AI/ML aplicado al RIC}

\subsubsection{Flujo de datos y entrenamiento de modelos}
\subsubsection{Inferencia y toma de decisiones}
\subsubsection{Closed-loop control basado en AI/ML}

\subsection{Plataformas experimentales y testbeds}

\subsubsection{OpenAirInterface y srsRAN}
\subsubsection{Colosseum y Arena}
\subsubsection{Plataformas PAWR y 5Genesis}


\section{Seguridad en Open RAN}

\subsection{Superficie de ataque de Open RAN}

\subsubsection{Impacto de la desagregación y multi-vendorismo}
\subsubsection{Virtualización, cloudificación y software de terceros}
\subsubsection{Nuevos vectores de ataque introducidos por las interfaces abiertas}

\subsection{Amenazas sobre los componentes e interfaces}

\subsubsection{Amenazas sobre O-RU, O-DU y O-CU}
\subsubsection{Amenazas sobre las interfaces E2, A1, O1 y O2}
\subsubsection{Amenazas sobre el Open Fronthaul}

\subsection{Amenazas asociadas al RIC y sus aplicaciones}

\subsubsection{xApps y rApps maliciosas o comprometidas}
\subsubsection{Ataques contra E2 Nodes y señalización}
\subsubsection{Conflictos y manipulación de decisiones de control}

\subsection{Ataques y riesgos de seguridad}

\subsubsection{Denial of Service y Distributed Denial of Service}
\subsubsection{Man-in-the-Middle y ataques sobre las comunicaciones}
\subsubsection{Spoofing, inyección de datos y manipulación de tráfico}
\subsubsection{Ataques relacionados con AI/ML}

\subsection{Privacidad y protección de información}

\subsubsection{Datos de usuarios y telemetría de la RAN}
\subsubsection{Riesgos derivados de la recopilación y exposición de datos}
\subsubsection{Privacidad en escenarios de inteligencia distribuida}


\section{Mecanismos y estrategias de seguridad para Open RAN}

\subsection{Seguridad de interfaces y comunicaciones}

\subsubsection{Autenticación y autorización}
\subsubsection{Protección de las comunicaciones mediante mecanismos criptográficos}
\subsubsection{Integridad y confidencialidad de la información}

\subsection{Control de acceso y confianza en el ecosistema O-RAN}

\subsubsection{Gestión de identidades y privilegios}
\subsubsection{Zero Trust aplicado a Open RAN}
\subsubsection{Verificación y aislamiento de aplicaciones}

\subsection{Seguridad de xApps y rApps}

\subsubsection{Validación y autenticación de aplicaciones}
\subsubsection{Gestión de permisos y privilegios}
\subsubsection{Detección de comportamiento malicioso}

\subsection{Detección y mitigación de ataques}

\subsubsection{Monitoreo de interfaces y tráfico}
\subsubsection{Detección basada en anomalías}
\subsubsection{Uso de AI/ML para seguridad}

\subsection{Estandarización y buenas prácticas}

\subsubsection{Esfuerzos de estandarización de seguridad}
\subsubsection{Interoperabilidad y pruebas de seguridad}
\subsubsection{Gestión del ciclo de vida y actualización de componentes}


\section{Relación entre RIC, inteligencia y seguridad}

\subsection{El RIC como habilitador de mecanismos de seguridad}

\subsubsection{Monitoreo inteligente de la RAN}
\subsubsection{Detección de anomalías en tiempo casi real}
\subsubsection{Respuesta y mitigación mediante closed-loop control}

\subsection{Seguridad de los mecanismos de inteligencia}

\subsubsection{Ataques contra modelos y datos AI/ML}
\subsubsection{Manipulación de modelos y decisiones}
\subsubsection{Seguridad del ciclo de vida AI/ML}

\subsection{Trade-offs entre inteligencia, apertura y seguridad}

\subsubsection{Beneficios de la apertura frente a la ampliación de la superficie de ataque}
\subsubsection{Flexibilidad multi-vendor frente a complejidad de integración}
\subsubsection{Automatización frente a riesgos asociados al control autónomo}


\section{Desafíos y oportunidades de investigación}

\subsection{Desafíos arquitectónicos y de interoperabilidad}

\subsubsection{Interoperabilidad multi-vendor}
\subsubsection{Evolución de interfaces y E2 Service Models}
\subsubsection{Escalabilidad y distribución cloud-edge}

\subsection{Desafíos de inteligencia y control}

\subsubsection{Latencia y restricciones temporales}
\subsubsection{Conflictos entre aplicaciones del RIC}
\subsubsection{Generalización y confiabilidad de modelos AI/ML}

\subsection{Desafíos de seguridad}

\subsubsection{Protección de componentes y aplicaciones de terceros}
\subsubsection{Seguridad de interfaces abiertas}
\subsubsection{Ataques contra AI/ML y control inteligente}

\subsection{Direcciones futuras}

\subsubsection{Evolución hacia redes 6G}
\subsubsection{Nuevos casos de uso y arquitecturas}
\subsubsection{Integración de inteligencia y seguridad por diseño}


\section{Análisis y discusión}

\subsection{Síntesis de la evolución hacia O-RAN inteligente y seguro}

\subsubsection{Relación entre apertura, desagregación e inteligencia}
\subsubsection{Relación entre RIC, automatización y optimización}
\subsubsection{Relación entre apertura, inteligencia y seguridad}

\subsection{Comparación de enfoques de implementación}

\subsubsection{Características de las principales implementaciones}
\subsubsection{Capacidades para desarrollo y experimentación}
\subsubsection{Limitaciones identificadas en la literatura}

\subsection{Análisis de los principales desafíos}

\subsubsection{Desafíos técnicos}
\subsubsection{Desafíos de seguridad}
\subsubsection{Desafíos de estandarización e interoperabilidad}


\section{Conclusiones}

\subsection{Principales hallazgos}

\subsubsection{Características fundamentales de Open RAN}
\subsubsection{Rol del RIC en la evolución de la RAN}
\subsubsection{Principales implicaciones de seguridad}

\subsection{Perspectivas futuras}

\subsubsection{Evolución del RIC y las arquitecturas O-RAN}
\subsubsection{Evolución de las estrategias de seguridad}
\subsubsection{Oportunidades de investigación en 5G-Advanced y 6G}

\end{document}
